\begin{definition}[Chain complex in a preadditive category] \label{definition:chain_complex_of_objects_in_a_preadditive_category}
Let $\mathcal{A}$ be a \CrefAndHyperrefIfExist{definition:additive_category_preadditive_category}{preadditive category} and let $I$ be a \CrefAndHyperrefIfExist{definition:total_order_linear_order_on_a_set}{totally ordered set} (typically $\mathbb{Z}$, but $I \subseteq \mathbb{Z}$ is also allowed). 
\begin{enumerate}
    \item A \hldef{chain complex} (or \hldef{homological (chain) complex}) $(K_\bullet, d_\bullet)$ in $\mathcal{A}$ indexed by $I$ is the homological convention for sequences with decreasing degrees. It consists of:
    \begin{itemize}
        \item Objects $\{ K_i \}_{i \in I}$ in $\mathcal{A}$, called the \hldef{terms in degree $i$},
        \item Morphisms $d_i: K_i \to K_{i-1}$ in $\mathcal{A}$, called the \hldef{boundary maps} or \hldef{differentials in degree $i$},
    \end{itemize}
    such that for every $i \in I$, $d_{i-1} \circ d_i = 0$. That is,
    $$ K_\bullet: \cdots \xrightarrow{d_{i+1}} K_i \xrightarrow{d_i} K_{i-1} \xrightarrow{d_{i-1}} K_{i-2} \xrightarrow{} \cdots $$
    with $d_{i-1}d_i = 0$ for all $i$. We typically use the notation \hl{$K_\bullet = (K_i, d_i)_{i \in I}$}.



    \item Dually, a \hldef{cochain complex} (or \hldef{cohomological complex}) $(K^\bullet, d^\bullet)$ in $\mathcal{A}$ follows the cohomological convention with increasing degrees. It consists of objects $\{ K^i \}_{i \in I}$ and \hldef{coboundary maps} $d^i: K^i \to K^{i+1}$ such that $d^{i+1} \circ d^i = 0$:
    $$ K^\bullet: \cdots \xrightarrow{d^{i-1}} K^i \xrightarrow{d^i} K^{i+1} \xrightarrow{d^{i+1}} K^{i+2} \xrightarrow{} \cdots $$
    We typically use the notation \hl{$K^\bullet = (K^i, d^i)_{i \in I}$}.

    \item Let $K_\bullet = (K_i, d_i^K)$ and $L_\bullet = (L_i, d_i^L)$ be \CrefAndHyperrefIfExist{definition:chain_complex_of_objects_in_a_preadditive_category}{chain complexes} in $\mathcal{A}$ indexed by the same set $I$. A \hldef{morphism of chain complexes} (or \hldef{chain map})
    $$ f_\bullet: K_\bullet \to L_\bullet $$
    consists of morphisms $f_i: K_i \to L_i$ for all $i \in I$, such that for every $i \in I$, the following diagram commutes:
    $$ \begin{array}{ccc} K_i & \xrightarrow{d_i^K} & K_{i-1} \\ \downarrow{f_i} && \downarrow{f_{i-1}} \\ L_i & \xrightarrow{d_i^L} & L_{i-1} \end{array} $$
    i.e., $d_i^L \circ f_i = f_{i-1} \circ d_i^K$. 



    A \hldef{morphism of cochain complexes} $f^\bullet: K^\bullet \to L^\bullet$ is defined similarly, satisfying the commutativity condition $d_L^i \circ f^i = f^{i+1} \circ d_K^i$.
\end{enumerate}

The collection of these objects and morphisms forms a category. Notation for these categories is as follows:
\begin{itemize}
    \item \hl{$\mathrm{Ch}(\mathcal{A})$} or \hl{$\mathbf{Ch}(\mathcal{A})$} is often used as a general term.
    \item To be explicit about the indexing convention, one uses \hl{$\mathrm{Ch}_\bullet(\mathcal{A})$} for chain complexes and \hl{$\mathrm{Ch}^\bullet(\mathcal{A})$} (or sometimes $\mathrm{CoCh}(\mathcal{A})$) for cochain complexes.
    \item The set of chain maps between two complexes is denoted by $\hlin{\operatorname{Hom}_{\mathrm{Ch}(\mathcal{A})}(K_\bullet, L_\bullet)}$; it is an abelian group under pointwise addition $(f+g)_i = f_i + g_i$.
\end{itemize}

Notations used to write chain complexes are \emph{homological conventions} whereas notations used to write cochain complexes are \emph{cohomological conventions}. In homological conventions, indices are written as subscripts and decrease when ``following the differential maps''. In cohomological conventions, indices are written as superscripts and incrase when ``following the diferential maps''. The theories of chain complexes and of cochain complexes are not genuinely different, however, as the conventions are essentially notational --- starting from a chain complex $(K_\bullet, d_\bullet)$, one can obtain a cochain complex $(K^\bullet, d^\bullet)$ simply by defining $K^i = K_{-i}$ and $d^i$ as $d_{-i}$ and vice versa. As such, cochain complexes may be simply called \hldef{chain complexes} when careful distinctions between the conventions are unnecessary. Homological and cohomological notational conventions persist throughout homological algebra, see e.g. \CrefAndHyperrefIfExist{definition:spectral_sequence_in_an_abelian_category}{spectral sequences}, \CrefAndHyperrefIfExist{definition:double_complex_of_objects_in_an_additive_category}{double complexes}, etc, but again the distinction between such conventions are merely notational.


\TextIfExists{definition:dg_category_over_a_ring}{
If $k$ is a \CrefAndHyperrefIfExist{definition:commutative_ring}{commutative ring} such that $\Hom_\calA(X,Y)$ is \CrefAndHyperrefIfExist{definition:category_enriched_in_a_monoidal_category}{enriched in} the category of \CrefAndHyperrefIfExist{definition:bimodule_over_rings_module_over_a_ring}{$k$-modules}, then $\mathrm{Ch}(\calA)$ \CrefAndHyperref{definition:category_of_chain_complexes_of_objects_in_an_additive_category_as_a_dg_category}{can be equipped with} the structure of a \CrefAndHyperrefIfExist{definition:dg_category_over_a_ring}{dg-category over $k$}.
}
\end{definition}

% \input{../_remarks/remark_cohomological_vs_homological_conventions.tex}
%\input{../_conventions/ogical_algebra_is_discussed_in_cohomological_terms.tex}
% \input{../_definitions/definition_morphism_of_chain_complexes_of_objects_in_an_additive_category.tex}

% See Also
% \input{../_concepts/proposition_category_of_chain_complexes_in_an_additive_category_is_additive.tex}